\documentclass[11pt]{article}

\usepackage{amsmath,amssymb,amsthm,mathtools}
\usepackage[margin=1in]{geometry}
\usepackage{longtable}
\usepackage{url}

\newtheorem{theorem}{Theorem}[section]
\newtheorem{lemma}[theorem]{Lemma}
\newtheorem{corollary}[theorem]{Corollary}
\newtheorem{remark}[theorem]{Remark}

\newcommand{\dd}{\,d}

\begin{document}

\section{Purpose and result}

For a finite graph $H$ and a graphon $W$ on an arbitrary probability space,
write
\[
 \Gamma_H(W)=t(H,W)-\Phi_H(p(W)),
 \qquad
 \Phi_H(p)=(1-p)^{v(H)}
 \chi_H\!\left(\frac1{1-p}\right).
\]
This note develops only the high-density and complete-multipartite
perturbation test.  It removes the exact-regularity hypothesis from the
previous regular-complement theorem.  Put $V=1-W$ and
\[
 q=\int_{\Omega^2}V,
 \qquad d(x)=\int_\Omega V(x,y)\dd y,
 \qquad \Delta=\lVert d\rVert_\infty.
\]
The sole analytic hypothesis is that the \emph{complement degree} is
uniformly small.  The complement itself may be arbitrary, nonregular,
non-step, and may take the value one.

Define the two universal nonregularity defects
\begin{equation}
 X_V=\int_Omega(d-q)^2
     =t(P_3,V)-q^2,
 \qquad
 Y_V=\int_{Omega^3}V_{12}V_{23}(1-V_{13})
     =t(P_3,V)-t(K_3,V).
 \label{eq:defects}
\end{equation}
Both are nonnegative.  The main theorem gives explicit graph constants
$c_H,\Lambda_H,N_3(H),\Omega_H$ such that
\begin{equation}
 \Gamma_H(1-V)\geq
 (c_H-\Lambda_H\Delta)X_V
 +(N_3(H)-\Omega_H\Delta)Y_V.
 \label{eq:main-preview}
\end{equation}
It follows that every one of the $29$ currently open rows is nonnegative in
an explicit, genuinely nonregular graphon neighbourhood of $p=1$.


\section{Tree densities with nonconstant degrees}

The first issue is that forest densities cease to equal $q^{e(F)}$ when
$d$ is nonconstant.  The following quantitative replacement is the needed
coercive estimate.

\begin{lemma}[Tree between the constant and the star]
\label{lem:tree-star}
Let $T$ be a tree with $e\geq1$ edges.  Then
\begin{equation}
 q^e\leq t(T,V)\leq\int_Omega d(x)^e\dd x.
 \label{eq:tree-star}
\end{equation}
Consequently, for $e\geq2$,
\begin{equation}
 0\leq t(T,V)-q^e
 \leq \binom e2\Delta^{e-2}X_V.
 \label{eq:tree-variance}
\end{equation}
For $e=1$ the difference is zero.
\end{lemma}

\begin{proof}
The result is immediate if $q=0$, so suppose $q>0$.  Let
\[
 \pi(\dd x)=\frac{d(x)}q\,\mu(\dd x),
 \qquad
 K(x,\dd y)=\frac{V(x,y)}{d(x)}\,\mu(\dd y)
\]
on the set $d(x)>0$, and define $K$ arbitrarily on the null set for $\pi$
where $d=0$.  Symmetry of $V$ shows that $K$ is reversible with stationary
law $\pi$.  Root $T$, sample the root from $\pi$, and independently sample
each child from $K$ conditional on its parent.  Every vertex variable $Z_u$
has marginal law $\pi$.  Comparing the joint density with the product
defining $t(T,V)$ gives the exact identity
\begin{equation}
 t(T,V)=q\,\mathbb E
 \prod_{u\in V(T)}d(Z_u)^{\deg_T(u)-1}.
 \label{eq:tree-markov}
\end{equation}
The exponents are nonnegative and sum to
$2e-(e+1)=e-1$.

For the lower bound, Jensen's inequality applied to the exponential gives
\[
 \mathbb E\prod_ud(Z_u)^{\deg_T(u)-1}
 \geq
 \exp\!\left((e-1)\int\log d\dd\pi\right).
\]
Moreover
\[
 \int\log(d/q)\dd\pi
 =\int \frac dq\log\frac dq\dd\mu\geq0,
\]
because $z\log z-z+1\geq0$ and $\int d/q=1$.  Thus the last display is at
least $q^{e-1}$.  Values where $d=0$ have zero $\pi$-mass; alternatively all
logarithms may first be truncated and then treated by monotone convergence.

For the upper bound, H\"older's inequality under the same joint law, using
exponent $(e-1)/(\deg_T(u)-1)$ whenever $\deg_T(u)>1$, gives
\[
 \mathbb E\prod_ud(Z_u)^{\deg_T(u)-1}
 \leq\prod_u
 \left(\mathbb E d(Z_u)^{e-1}\right)^{(\deg_T(u)-1)/(e-1)}
 =\int d^{e-1}\dd\pi.
\]
After multiplication by $q$, this is $\int d^e$.  This proves
\eqref{eq:tree-star} without a finite-graph approximation.

For $z,q\in[0,\Delta]$, Taylor's theorem and the upper bound
$e(e-1)\Delta^{e-2}$ on the second derivative of $z^e$ give
\[
 z^e\leq q^e+eq^{e-1}(z-q)
       +\binom e2\Delta^{e-2}(z-q)^2.
\]
Integrating, the linear term vanishes and the quadratic integral is $X_V$.
Combine this with \eqref{eq:tree-star}.
\end{proof}


\section{Every long failed closure is controlled by a triangle failure}

For $\ell\geq3$, define the open-cycle defect
\[
 Z_\ell(V)=\int_{\Omega^\ell}
 \left(\prod_{i=0}^{\ell-2}V(x_i,x_{i+1})\right)
 (1-V(x_0,x_{\ell-1}))\dd\mu^\ell.
\]
Thus $Z_3(V)=Y_V$.

\begin{lemma}[Path-to-triangle telescoping]
\label{lem:path-triangle}
For every $\ell\geq3$,
\begin{equation}
 0\leq Z_\ell(V)\leq(\ell-2)\Delta^{\ell-3}Y_V.
 \label{eq:path-triangle}
\end{equation}
\end{lemma}

\begin{proof}
For any $a_1,\ldots,a_s\in[0,1]$,
\begin{equation}
 a_1(1-a_s)\leq\sum_{i=1}^{s-1}a_i(1-a_{i+1}).
 \label{eq:scalar-chain}
\end{equation}
Indeed, the two-step inequality
\[
 a(1-c)\leq a(1-b)+b(1-c)
\]
has remainder $b(1-a-c)+ac$.  If $a+c\leq1$ it is visibly nonnegative;
if $a+c\geq1$, its minimum over $b\in[0,1]$ is
$(1-a)(1-c)$.  Iteration proves \eqref{eq:scalar-chain}.

Apply it pointwise with $a_i=V(x_0,x_i)$ for
$1\leq i\leq\ell-1$, and multiply by the remaining consecutive path-edge
factors.  Each of the $\ell-2$ resulting terms contains precisely the
triangle-failure factor
\[
 V(x_0,x_i)V(x_i,x_{i+1})(1-V(x_0,x_{i+1})).
\]
The other $\ell-3$ vertices lie in two tree branches attached to this
triple.  Integrating them successively from their leaves contributes at most
$\Delta^{\ell-3}$.  The remaining triple integral is $Y_V$.
\end{proof}

Let $C$ be a connected graph.  For a spanning tree $T\subseteq C$, put
\[
 \kappa_T(C)=\sum_{uv\in E(C)\setminus E(T)}
 \bigl(\operatorname{dist}_T(u,v)-1\bigr),
 \qquad
 \kappa(C)=\min_T\kappa_T(C).
\]
For a tree, $\kappa(C)=0$.

\begin{lemma}[Connected-component defect]
\label{lem:connected}
If $C$ is connected on $v\geq3$ vertices and $r=v-1$, then
\begin{equation}
 \left|t(C,V)-q^r\right|
 \leq \binom r2\Delta^{r-2}X_V
      +\kappa(C)\Delta^{r-2}Y_V.
 \label{eq:connected}
\end{equation}
\end{lemma}

\begin{proof}
Choose a spanning tree $T$ attaining $\kappa(C)$.  Since $0\leq V\leq1$,
$t(C,V)\leq t(T,V)$.  Also
\[
 t(T,V)-t(C,V)
 \leq\sum_{uv\in E(C)\setminus E(T)}
 \int\left(\prod_{ab\in E(T)}V_{ab}\right)(1-V_{uv}).
\]
For a fixed extra edge $uv$, its path in $T$ closes a cycle of length
$\operatorname{dist}_T(u,v)+1$.  Integrate every branch off that path by the
degree bound $d\leq\Delta$, and apply Lemma~\ref{lem:path-triangle}.  The term
is at most
\[
 \bigl(\operatorname{dist}_T(u,v)-1\bigr)
 \Delta^{v-3}Y_V.
\]
Summing gives
$t(T,V)-t(C,V)\leq\kappa(C)\Delta^{r-2}Y_V$.
Lemma~\ref{lem:tree-star} bounds
$t(T,V)-q^r$ by the first term of \eqref{eq:connected}.  Both $t(C,V)$ and
$q^r$ are at most $t(T,V)$, so their absolute difference is at most the sum
of these two upper defects.
\end{proof}


\section{Graphic-rank expansion and the coercive theorem}

For $A\subseteq E(H)$, let $H_A=(V(H),A)$ and
\[
 r_H(A)=v(H)-c(H_A).
\]
For every nontrivial connected component $C$ of $H_A$, put
\[
 \lambda(C)=\binom{|V(C)|-1}{2}.
\]
Define
\[
 \lambda(A)=\sum_C\lambda(C),
 \qquad
 \omega(A)=\sum_C\kappa(C),
\]
where isolated vertices contribute zero, and finally
\begin{equation}
 \Lambda_H=\sum_{A:r_H(A)\geq3}\lambda(A),
 \qquad
 \Omega_H=\sum_{A:r_H(A)\geq3}\omega(A).
 \label{eq:global-constants}
\end{equation}

\begin{lemma}[Arbitrary edge-subset defect]
\label{lem:subset}
If $R=r_H(A)\geq2$, then
\begin{equation}
 \left|t(H_A,V)-q^R\right|
 \leq\lambda(A)\Delta^{R-2}X_V
     +\omega(A)\Delta^{R-2}Y_V.
 \label{eq:subset}
\end{equation}
\end{lemma}

\begin{proof}
Both $t(H_A,V)$ and $q^R$ factor over the nontrivial components.  A connected
component of rank $r$ has density at most that of a spanning tree, hence at
most $\Delta^r$ by leaf integration; also $q^r\leq\Delta^r$.  Telescope the
difference of the two component products.  Lemma~\ref{lem:connected} bounds
the selected component, and all unchanged factors contribute at most
$\Delta^{R-r}$.  Summing over components proves \eqref{eq:subset}.  Components
of order two are exact because $t(K_2,V)=q$.
\end{proof}

Whitney's edge-subset expansion and inclusion--exclusion give, for $q>0$,
\begin{equation}
 \Gamma_H(1-V)=\sum_{A\subseteq E(H)}(-1)^{|A|}
 \bigl(t(H_A,V)-q^{r_H(A)}\bigr).
 \label{eq:rank-expansion}
\end{equation}
Indeed,
$q^{v(H)}\chi_H(1/q)=\sum_A(-1)^{|A|}q^{r_H(A)}$.
The identity is polynomial in $q$, so its $q=0$ endpoint is obtained directly
or by continuity.

Let $N_3(H)$ be the number of triangles and let
\[
 c_H=\sum_{x\in V(H)}\binom{\deg_H(x)}2-N_3(H).
\]

\begin{theorem}[Maximum-degree complement stability]
\label{thm:main}
For every finite simple graph $H$ and every graphon $V$,
\begin{equation}
 \Gamma_H(1-V)\geq
 (c_H-\Lambda_H\Delta)X_V
 +(N_3(H)-\Omega_H\Delta)Y_V.
 \label{eq:main}
\end{equation}
Consequently, if all four graph constants are positive and
\begin{equation}
 \Delta<\eta_H:=
 \min\left\{\frac{c_H}{\Lambda_H},
             \frac{N_3(H)}{\Omega_H}\right\},
 \label{eq:radius}
\end{equation}
then $\Gamma_H(1-V)\geq0$.  Equality can occur only when $V=0$ almost
everywhere or when $1-V$ is a balanced Turan graphon.
\end{theorem}

\begin{proof}
The rank-zero and rank-one summands in \eqref{eq:rank-expansion} vanish.  At
rank two, every pair of disjoint edges also vanishes.  Every adjacent pair
contributes $X_V$.  A triangle contributes
\[
 -\bigl(t(K_3,V)-q^2\bigr)=-X_V+Y_V.
\]
There are $\sum_x\binom{\deg_H(x)}2$ adjacent pairs and $N_3(H)$ triangles,
so the complete rank-two contribution is
$c_HX_V+N_3(H)Y_V$.

For $R\geq3$, Lemma~\ref{lem:subset} and $0\leq\Delta\leq1$ give
\[
 \left|t(H_A,V)-q^R\right|
 \leq\lambda(A)\Delta X_V+\omega(A)\Delta Y_V.
\]
Discard every sign and sum.  This proves \eqref{eq:main}.

Under the strict condition \eqref{eq:radius}, equality forces $X_V=Y_V=0$.
If $X_V=0$, then $d=q$ almost everywhere.  If also $Y_V=0$, then
\[
 V(x,y)V(y,z)>0\quad\Longrightarrow\quad V(x,z)=1
\]
for almost every triple.  Here is the support-class argument explicitly.
For typical $x$, let $N_x=\{z:V(x,z)>0\}$.  Then
$\mu(N_x)\geq d(x)=q$.  For almost every $y\in N_x$, transitivity applied to
$y,x,z$ gives $V(y,z)=1$ for almost every $z\in N_x$.  Hence
\[
 q=d(y)\geq\mu(N_x)\geq q.
\]
All inequalities are equalities, so $\mu(N_x)=q$, $V=1$ almost everywhere
on $N_x^2$, and the degree identity forces $V=0$ from $N_x$ to its
complement.  Therefore two typical positive-support neighbourhoods are
either equal modulo null sets or disjoint.  They partition the probability
space into blocks of mass $q$.  Thus their number is $r=1/q$, an integer,
and $1-V=T_r$ up to null sets.  The case $q=0$ is $V=0$ and $W=1$.
\end{proof}


\section{The 29 exact certificates}

The graphs use the NetworkX Atlas labelling and vertex set
$\{0,\ldots,n-1\}$.  The following data identify them independently of the
drawings in the catalogue.  The final column points to the exact chromatic
polynomial and target type in the next table.

\scriptsize
\begin{longtable}{r@{\quad}lrrr@{\quad}p{8.2cm}}
Atlas&graph6&$n$&$m$&type&edge set\\ \hline
43&\verb|Dlc|&5&6&1&01,03,04,12,23,34\\
118&\verb|Eht?|&6&7&2&01,04,12,14,15,23,34\\
122&\verb|Eld?|&6&7&2&01,03,04,12,15,23,34\\
124&\verb|Exd?|&6&7&2&01,02,04,12,15,23,34\\
126&\verb|ERUO|&6&7&2&02,05,13,14,23,34,35\\
127&\verb|EZEG|&6&7&3&02,05,12,13,23,34,45\\
130&\verb|E{CW|&6&7&4&01,02,03,12,34,35,45\\
137&\verb|EjsG|&6&8&5&01,04,12,13,14,23,34,45\\
139&\verb|Eju?|&6&8&5&01,04,05,12,13,14,23,34\\
145&\verb|Exe_|&6&8&6&01,02,04,05,12,23,25,34\\
147&\verb|EhMg|&6&8&6&01,05,12,23,24,25,34,45\\
148&\verb|EyUG|&6&8&6&01,02,05,12,13,14,34,45\\
151&\verb|EhdW|&6&8&7&01,04,12,15,23,34,35,45\\
152&\verb|EhNG|&6&8&6&01,05,12,15,23,24,34,45\\
153&\verb|Ehf_|&6&8&8&01,04,05,12,15,23,25,34\\
157&\verb|E~`_|&6&9&9&01,02,03,04,12,13,15,23,25\\
160&\verb|E~@g|&6&9&9&01,02,03,12,13,15,23,25,45\\
168&\verb|E^MG|&6&9&10&02,03,05,12,13,23,24,34,45\\
169&\verb|Exf_|&6&9&11&01,02,04,05,12,15,23,25,34\\
171&\verb|Ehew|&6&9&12&01,04,05,12,23,25,34,35,45\\
174&\verb|EtTg|&6&9&13&01,02,03,14,15,23,25,34,45\\
178&\verb|En}?|&6&10&14&01,03,04,05,12,13,14,23,24,34\\
181&\verb|E^mG|&6&10&15&02,03,04,05,12,13,23,24,34,45\\
185&\verb|Exv_|&6&10&16&01,02,04,05,12,14,15,23,25,34\\
188&\verb|EzNG|&6&10&17&01,02,05,12,13,15,23,24,34,45\\
194&\verb|E~wW|&6&11&18&01,02,03,04,12,13,14,23,24,35,45\\
196&\verb|ER~g|&6&11&18&02,04,05,13,14,15,23,24,25,34,45\\
199&\verb|El^g|&6&11&19&01,03,05,12,14,15,23,24,25,34,45\\
203&\verb|E~^G|&6&12&20&01,02,03,05,12,13,14,15,23,24,34,45
\end{longtable}
\normalsize

All rows of types $1$--$8$, $10$, $12$, $13$, and $17$ have chromatic
number three; types $9$, $11$, $14$--$16$, and $18$--$20$ have chromatic
number four.  This follows both from the first nonzero positive integer of
the displayed polynomial and from the explicit triangle or $K_4$ visible in
the edge list.  Deletion--contraction gives the following exact expressions.

\scriptsize
\begin{longtable}{r@{\quad}p{7.1cm}p{7.1cm}}
type&$\chi_H(x)$&$\Phi_H(p)$\\ \hline
1&$x(x-1)(x-2)(x^2-3x+3)$&$p(2p-1)(3p^2-3p+1)$\\
2&$x(x-1)^2(x-2)(x^2-3x+3)$&$p^2(2p-1)(3p^2-3p+1)$\\
3&$x(x-1)(x-2)^2(x^2-2x+2)$&$p(2p-1)^2(2p^2-2p+1)$\\
4&$x(x-1)^3(x-2)^2$&$p^3(2p-1)^2$\\
5&$x(x-1)^2(x-2)^3$&$p^2(2p-1)^3$\\
6&$x(x-1)(x-2)^2(x^2-3x+3)$&$p(2p-1)^2(3p^2-3p+1)$\\
7&$x(x-1)(x-2)^2(x^2-3x+4)$&$p(2p-1)^2(4p^2-5p+2)$\\
8&$x(x-1)(x-2)(x^3-5x^2+9x-7)$&$p(2p-1)(7p^3-12p^2+8p-2)$\\
9&$x(x-1)^2(x-2)^2(x-3)$&$p^2(2p-1)^2(3p-2)$\\
10&$x(x-1)(x-2)^2(x^2-4x+5)$&$p(2p-1)^2(5p^2-6p+2)$\\
11&$x(x-1)(x-2)(x-3)(x^2-3x+3)$&$p(2p-1)(3p-2)(3p^2-3p+1)$\\
12&$x(x-1)(x-2)(x^3-6x^2+13x-11)$&$p(2p-1)(11p^3-20p^2+13p-3)$\\
13&$x(x-1)(x-2)(x^3-6x^2+14x-13)$&$p(2p-1)(13p^3-25p^2+17p-4)$\\
14&$x(x-1)^2(x-2)(x-3)^2$&$p^2(2p-1)(3p-2)^2$\\
15&$x(x-1)(x-2)^3(x-3)$&$p(2p-1)^3(3p-2)$\\
16&$x(x-1)(x-2)(x-3)(x^2-4x+5)$&$p(2p-1)(3p-2)(5p^2-6p+2)$\\
17&$x(x-1)(x-2)(x^3-7x^2+18x-17)$&$p(2p-1)(17p^3-33p^2+22p-5)$\\
18&$x(x-1)(x-2)(x-3)(x^2-5x+7)$&$p(2p-1)(3p-2)(7p^2-9p+3)$\\
19&$x(x-1)(x-2)(x-3)(x^2-5x+8)$&$p(2p-1)(3p-2)(8p^2-11p+4)$\\
20&$x(x-1)(x-2)(x-3)(x^2-6x+10)$&$p(2p-1)(3p-2)(10p^2-14p+5)$
\end{longtable}
\normalsize

Exact enumeration of every edge subset and every spanning tree in its
nontrivial components gives the following theorem constants.  No decimal
rounding is involved.

\scriptsize
\begin{longtable}{r|rrrr|c@{\qquad}r|rrrr|c}
$H$&$c_H$&$\Lambda_H$&$N_3$&$\Omega_H$&$\eta_H$&
$H$&$c_H$&$\Lambda_H$&$N_3$&$\Omega_H$&$\eta_H$\\ \hline
43&8&161&1&18&$8/161$&118&11&481&1&37&$11/481$\\
122&10&461&1&37&$10/461$&124&10&464&1&37&$5/232$\\
126&10&454&1&31&$5/227$&127&9&454&1&31&$9/454$\\
130&8&376&2&30&$1/47$&137&14&1173&3&131&$14/1173$\\
139&13&1134&3&131&$13/1134$&145&14&1211&2&131&$2/173$\\
147&13&1184&2&126&$13/1184$&148&13&1207&2&122&$13/1207$\\
151&13&1237&1&138&$1/138$&152&12&1159&2&122&$12/1159$\\
153&12&1174&2&128&$6/587$&157&17&2725&5&423&$17/2725$\\
160&16&2638&5&423&$8/1319$&168&17&2893&3&421&$17/2893$\\
169&16&2778&4&406&$8/1389$&171&16&2903&3&413&$16/2903$\\
174&16&2972&2&436&$1/218$&178&20&6138&7&1213&$10/3069$\\
181&20&6482&6&1123&$10/3241$&185&20&6598&5&1193&$10/3299$\\
188&20&6775&4&1215&$4/1355$&194&24&14648&7&3229&$3/1831$\\
196&24&14852&7&3091&$6/3713$&199&24&14992&6&3276&$3/1874$\\
203&28&32316&9&8052&$7/8079$&&&&&&
\end{longtable}
\normalsize

\begin{corollary}[All current open rows]
For every graph $H$ in the tables and every graphon $W$, if
\[
 \operatorname*{ess\,sup}_{x}\int(1-W(x,y))\dd y<\eta_H,
\]
then $t(H,W)\geq\Phi_H(p(W))$.  These densities lie inside the required
range: every $\eta_H\leq8/161<1/3$, whereas
$1/(\chi(H)-1)$ is $1/2$ or $1/3$.
\end{corollary}

At $p=1$, the complement has $q=0$ and hence $V=0$ almost everywhere.
Thus $W=1$, $t(H,W)=1$, and the polynomial continuation of the monic
chromatic polynomial gives $\Phi_H(1)=1$.  This handles the endpoint without
dividing by $q$.


\section{Hidden-assumption audit and limitations}

All integrations above concern bounded measurable functions, so Tonelli and
Fubini apply.  The tree-indexed Markov law is constructed directly on the
original probability space; no finite graph, step approximation, compactness
argument, or regularity lemma is used.  The proof never assumes that $d$ is
constant.  The only pointwise hypothesis is the explicitly displayed bound
$d\leq\Delta$ almost everywhere.  There is no conditioning on an event of
possibly zero mass: the case $q=0$ is separated before the Markov kernel is
defined.

This theorem by itself does not cover all graphons of small \emph{average}
complement density: a small exceptional set may still have large complement
degree.  The companion
\path{notes/arbitrary_graphon_high_density_stability.tex} applies the present
theorem to the low-degree residual after splitting off exactly that
high-degree set; its quantitative complete-core margin absorbs the mixed
target error.  The combination closes the high-density endpoint for every
graphon.  It still does not prove any catalogue row on its complete density
interval, so no status is changed.  The constants here are deliberately
obtained by taking absolute values of every graphic-rank-at-least-three
summand, and further cancellation could enlarge the final endpoint radii.

Run
\begin{verbatim}
python codes/principled_negative_tests.py
\end{verbatim}
from the repository root.  The line headed
``maximum-degree complement graphon high-density theorem'' reconstructs
$c_H,\Lambda_H,N_3(H),\Omega_H$, and $\eta_H$ for all $29$ rows from their
edge sets.  It enumerates every relevant edge subset and every spanning tree
using exact integer and rational arithmetic.

For formal verification, the reusable analytic layer should contain
Lemma~\ref{lem:tree-star}, Lemma~\ref{lem:path-triangle}, the component-product
telescoping lemma, and Theorem~\ref{thm:main}.  The finite layer should then
certify the four integer constants for the $29$ Atlas modules.  Since this is
a restricted high-density theorem and not a full classification, none of the
Atlas status modules should be changed.

\end{document}
